\chapter*{致谢}\addcontentsline{toc}{chapter}{致谢}
本次《软件工程综合实践》课程设计从需求分析、系统设计推进到四端实现、联调与答辩，离不开李罡老师、王立老师和张怡老师的指导。在此，我们向三位老师表示衷心感谢。老师们对需求、协作方式和技术方案的追问，使我们没有停留在“功能做出来了”的满足之中，而是进一步思考：这个功能究竟解决什么问题，我们能否解释它的实现，以及它在更复杂的使用条件下是否仍然成立。

尤其感谢李罡老师在项目初期对我们 SRS 提出的意见。刚开始时，我们更容易把注意力放在页面、技术和新增功能上，希望通过更多的模块呈现项目成果，却没有充分说明用户为什么需要这些功能、不同角色如何协作，以及异常情况应当如何处理。老师的意见让我们明白，产品要先面向需求，技术实现应当建立在对需求的理解之上。

带着这一认识，我们开始重新完善 SRS，梳理顾客、商家、骑手和管理员的目标，明确订单与配送的关系，补充取消、库存不足、权限限制和外部服务失败等情形。文档的完善也促使我们重新审视代码：一个状态名称是否清楚，一项优惠规则是否有统一口径，一个接口是否能被验证，都与需求是否充分表达有关。我们由此逐渐理解，课程设计中的需求文档并不是最后提交时补上的说明，而是帮助整个团队建立共同理解的基础。

感谢李罡老师、王立老师和张怡老师在答辩中围绕分工协作提出的问题。当老师们询问我们认为按端开发更好，还是全栈开发更好时，我们开始意识到，分工不能只以“每个人做了几个页面”衡量。顾客下单、商家出餐和骑手配送属于同一条业务链路，任何一端都无法独立证明整个系统已经完成。按端主责有助于保持角色体验，而围绕业务用例贯通前后端、数据与测试，能够让成员真正理解自己交付的结果。对我们而言，更重要的是在责任明确的同时保持共同讨论、及时联调与相互评审。

老师们关于 AI 点餐的追问也使我们受益良多。对于“回答一个问题调用三次是否过多，并发量增加后会发生什么、如何解决”的问题，我们最初更关注 AI 是否能够理解输入、返回答案，却对等待时间、接口限额、重试放大和资源占用考虑得不够充分。通过重新检查调用链路，我们认识到，有些问题可以直接由规则与真实业务数据回答；需要使用模型时，也应当说明每次调用的必要性，并通过调用预算、限流、超时、降级和测试控制影响。技术的价值不仅在于能够展示新的能力，也在于能够以合理的成本和清晰的边界服务用户。

我们也进一步认识到 AI 辅助开发的责任边界。工具可以帮助我们理解报错、整理材料和形成方案，但需求取舍、架构判断和测试结果仍需要由我们自己理解和确认。只有能够解释实现的原因、识别它的不足，并在失败场景中验证它，工具的帮助才真正转化为我们的能力。

感谢各位小组成员在本次实践中的共同投入。范文丽同学负责前端与看板，李承文同学负责后端与红包积分，赵紫怡同学负责数据库与 AI，仲禹潼同学负责测试与高德地图。不同方向的工作需要在同一业务流程中衔接，让我们体会到清楚的接口约定和及时沟通的重要性。需求讨论、接口约定、代码修改、问题排查和文档整理都需要持续沟通。项目中的分歧和反复也提醒我们，合作不只是把任务划开，更是让每个人的工作能够被其他成员理解，并最终组成一致的系统。

回顾这次实践，我们仍有需求与实现同步不够及时、测试证据不够完整和工程经验不足之处。感谢三位老师帮助我们看见这些问题，也让我们理解了课程设计的意义：从真实需求出发，在约束条件下做出有理由的选择，通过协作与验证形成成果，并对不足保持清楚的认识。我们将把这些收获带到后续学习与项目实践中，继续完善自己的需求分析、系统设计和工程协作能力。

\vspace{1em}\begin{flushright}轻量级外卖服务平台项目组\\2026 年 9 月\end{flushright}
